\chapter{Formats de musique}
\soft{} reconnaît \lhunaModulesCount{} formats de musique. Nous avons généré la table ci-dessous à partir de
l'application elle-même ; elle est donc exacte pour la version décrite.

Derrière ce nombre se cachent quatre façons très différentes de conserver une musique. Savoir à laquelle appartient
un fichier vous dit ce que vous pouvez en attendre : un module de tracker se rejoue toujours à l'identique, un
enregistrement de registres ne contient que ce qui a été capté, un rip exécute le vrai programme de la machine et
réagit donc à l'émulation dans ses moindres détails.

\begin{description}
  \item[Tracker module] partition et instruments, rejoués par un lecteur fidèle à celui d'origine ;
  \item[Register stream] suite de valeurs écrites dans les registres de la puce, enregistrées depuis la machine ;
  \item[Memory dump] programme de lecture d'origine exécuté par un processeur émulé (Z80, 6502\dots) ;
  \item[Multichip] plusieurs puces jouées ensemble.
\end{description}


{\small\lhunaFormatsTable}

\section{Conteneurs}
\soft{} ouvre aussi \lhunaContainersCount{} formats de conteneurs : archives, compresseurs, images disque,
instantanés de machine et décompilateurs de lecteurs.

{\small\lhunaContainersTable}
